

\selectlanguage{ngerman}%
\frenchspacing%
\begin{ledgroupsized}[r]{120mm}%
\footnotesize%
\pstart%
\noindent\textbf{Überlieferung:}%
\pend%
\end{ledgroupsized}%
\begin{ledgroupsized}[r]{114mm}%
\footnotesize%
\pstart%
\parindent -6mm
\makebox[6mm][l]{\textit{L}}%
Konzept:
LH~XXXVII~4 Bl.\;169--170.
Ein Bogen 2\textsuperscript{o};
Wasserzeichen in Bl.~\href{\Scan{LH_37_04_169wz}}{169} (Papier aus Rom).
Vier Seiten.
Textfolge: Bl.~169\:r\textsuperscript{o}, 169\:v\textsuperscript{o}, 170\:r\textsuperscript{o}, 170\:v\textsuperscript{o}, 169\:r\textsuperscript{o}.
Die uneinheitliche Schreibweise des Namens Guericke wird beibehalten.
\pend%
\end{ledgroupsized}
\selectlanguage{latin}%
\frenchspacing%
\vspace{8mm}%
\normalsize%
\pstart%
\noindent%
\centering%
\href{\Scan{LH_37_04_169r}}%
{\lbrack169\:r\textsuperscript{o}\rbrack}\
\edlabel{LH_37_04_169-170_08a}%
\hfill
Problema
\hfill
\edtext{}{%
\xxref{LH_37_04_169-170_08a}{LH_37_04_169-170_08b}%
\lemma{Problema}%
\Bfootnote{%
\textit{(1)}~Indices %
\textit{(2)}~Indicem
\textit{(a)}~exactum
\textit{(b)}~accuratum
\textit{(3)}~Exactum Indicem~\textit{L}}}
\linebreak
\phantom{\lbrack169\:r\textsuperscript{o}\rbrack}
Exactum Indicem%
\edlabel{LH_37_04_169-170_08b}
\linebreak
\phantom{\lbrack169\:r\textsuperscript{o}\rbrack}
densitatum\indexS{densitas aeris} et raritatum Aeris\indexS{raritas aeris}
\edlabel{LH_37_04_169-170_15a}
designare seu Thermometrum\indexS{thermometrum} exhibere ad
\linebreak
\phantom{\lbrack169\:r\textsuperscript{o}\rbrack}
gradus caloris\indexS{gradus caloris} vel frigoris\indexS{gradus frigoris} designandos accurate divisurum
\edlabel{LH_37_04_169-170_15b}
\edtext{}{%
\xxref{LH_37_04_169-170_15a}{LH_37_04_169-170_15b}
\lemma{designare}%
\Bfootnote{%
\textit{(1)}~, quod pro aere est %
\textit{(2)}~seu Thermometrum %
\textit{(a)}~exactum %
\textit{(b)}~accuratum %
\textit{(c)}~exhibere %
\textit{(aa)}~et %
\textit{(bb)}~convenientem ejus divisionem instituere %
\textit{(cc)}~accurate divis %
\textit{(dd)}~ad %
\textit{(aaa)}~aequales %
\textit{(bbb)}~gradus caloris %
\textit{(aaaa)}~accurate divisum %
\textit{(bbbb)}~vel frigoris designandos accurate divisurum~\textit{L}}}
\pend
\nobreak
\vspace{0.5\baselineskip}
\pstart%

\noindent%
Sit materiae aptae, ut aeris\indexS{aer} communis, certa
quantitas 
\edtext{datae densitatis,\indexS{densitas}}{\lemma{datae}\Bfootnote{%
\textit{(1)}~quantitatis %
\textit{(2)}~densitatis~\textit{L}}}
in datae magnitudinis
spatio inclusa. Jam qua ratione
augetur spatium, quod eadem moles implet,
in eadem
\edtext{ratione minuitur}{\lemma{ratione}\Bfootnote{%
\textit{(1)}~augetur %
\textit{(2)}~minuitur~\textit{L}}}
densitas\indexS{densitas} %
\edlabel{LH_37_04_169-170_21a}%
(per %
\edtext{prop.\,6 %
\edtext{cap. de quantitate%
\edtext{materiae); sufficit}{\lemma{materiae)}
\Bfootnote{%
\textit{(1)}~neque opus est ponderis coercentis rationem haberi %
\textit{(2)}~; sufficit~\textit{L}}}%
}{\lemma{(per}\Bfootnote{
\textit{(1)}~...) %
\textit{(2)}~prop.\,6 cap. de quantitate materiae)~\textit{L}}}
}{\lemma{prop.\,6 \lbrack...\rbrack\ materiae}%
\Cfootnote{\textsc{G.~W. Leibniz}, \textit{Dynamica de potentia et legibus}, %
pars\,I, sectio\,I, cap.\,2, prop.\,6 (%
\textit{LMG}\cite{01043}~VI, S.\,300\cite{01354}; %
\textit{CFP}~I, S.\,35\:\lbrack A\rbrack, II.1, S.\,42\textendash45\:\lbrack C\rbrack\cite{01482}).%
\hfill
}}%
\edlabel{LH_37_04_169-170_21b}
ergo quantitates spatiis reciprocas
numeris designari, ut habeatur index
densitatum\indexS{index densitatum} et raritatum\indexS{index raritatum}. Hinc si%
\edtext{spatium est ubique}{\lemma{spatium}
\Bfootnote{%
\textit{(1)}~semper %
\textit{(2)}~est ubique~\textit{L}}}
uniforme\lbrack,\rbrack\ dividendum est \textls{progressione harmonica}.
Ut si tubus sit vitreus\indexS{tubus vitreus} \textit{LMN} clausus %
in \textit{L}\lbrack,\rbrack\ apertus in \textit{N},
\edtext{(et aequalis amplitudinis in}{\lemma{et}\Bfootnote{%
\textit{(1)}~in parte %
\textit{(2)}~aequalis amplitudinis in~\textit{L}}}
\textit{L\ind{1}} parte
ipsius %
\edtext{\textit{LM}); tunc }{\lemma{}\Bfootnote{\textit{LM}) \ \textbar \ aequalis amplitudinis \textit{ gestr.}\ \textbar\ ; tunc \textit{ L}}}%
ut Aer\indexS{aer} %
in \textit{L1} sit densitatis\indexS{densitas} ut 1, in \textit{L2} ut 2, in \textit{L3} ut 3
et ita porro, debent spatia \textit{L1}, \textit{L2}, \textit{L3} et caetera
esse progressionis harmonicae, et secundum \textit{L2}
dimidium esset primi \textit{L1}, unde habentur et reliqua. %
Ita enim spatia erunt densitatibus\indexS{densitas} reciproce %
proportionalia. %
\edtext{Quae divisio sufficeret in thermometris\indexS{thermometrum} vulgaribus, apertis,
uti ponderis in liquore non habetur ratio\indexS{ratio ponderum}. Sed notum est %
haec thermometra\indexS{thermometrum} %
etiam a pondere %
aeris\indexS{pondus aeris} pati, nec %
exacta esse, nisi adsit %
praeterea barometrum, %
cujus comparatione corrigatur aestimatio.
}{\lemma{}\Bfootnote{Quae \lbrack...\rbrack\ aestimatio.~\textit{erg.~L}}}
\pend
\begin{edfigure}[1.5em]%
\hfill%
\hspace*{5mm}%
\includegraphics[scale=1]{%
gesamttex/edit_VIII-4/images/LH_37_04_169-170_d1_169r.pdf}%
\hfill%
\includegraphics[scale=1]{%
gesamttex/edit_VIII-4/images/LH_37_04_169-170_d2_169r.pdf}%
\hfill%
\edcaption[1em]{%
\lbrack\textit{Fig.~1}\rbrack%
\hspace*{41mm}%
\lbrack\textit{Fig.2}\rbrack%
\hspace*{5mm}}%
\end{edfigure}%
\pstart
At si non %
\edtext{tam}{\lemma{}\Bfootnote{tam~\textit{erg.~L}}}
densitatum\indexS{densitas} %
\edtext{vel raritatum}{\lemma{}\Bfootnote{vel raritatum~\textit{erg.~L}}}
gradus\indexS{gradus densitatum}\indexS{gradus raritatum} %
\edtext{in aere\indexS{aer} incluso}{\lemma{}\Bfootnote{in aere incluso~\textit{erg.~L}}}%
\edtext{quam caloris\indexS{calor} et frigoris%
\edtext{hoc est}{\lemma{}
\Afootnote{%
\textit{Am Rand, gestrichen:} \footnotesize{Hoc est si desideretur \textls{Thermometrum\indexS{thermometrum}}} \lbrack\textit{Text bricht ab.}\rbrack
}}
densitatis\indexS{densitas} et raritatis\indexS{raritas} in aere\indexS{aer} ambiente quaerantur\lbrack,\rbrack\ id est si desideretur thermometrum\indexS{thermometrum} clausum accurate divisum, sic rem instituemus in exemplo. Sit ex vitro instrumentum\indexS{instrumentum} \textit{MPQRSTVXZ}, cujus pars \textit{MPQR}
}{\lemma{quam}
\Bfootnote{%
\textit{(1)}~calorum quaerantur, sic rem instituemus in exemplo. Sit \bars{ex vitro~\textit{streicht Hrsg.}} tubus recurvus cum globo cavo \textit{OPQRSXVZ}, si \textit{C} %
\textit{(2)}~caloris %
\textit{(a)}~et frigoris (hoc est densitatis et raritatis) in ambiente aere %
\textit{(b)}~et frigoris hoc est densitatis et raritatis in aere ambiente quaerantur %
\textit{(aa)}~hoc %
\textit{(bb)}~id est si desideretur thermometrum clausum %
\textit{(aaa)}~instituemus \bars{in~\textit{streicht Hrsg.}} %
\textit{(bbb)}~accurate divisum, sic rem instituemus in exemplo. Sit ex vitro %
\textit{(aaaa)}~machina\indexS{machina}%
\textit{(bbbb)}~instrumentum
\textit{(aaaaa)}~\textit{OPQRSTVXZ} %
\textit{(bbbbb)}~\textit{MPQRSTVXZ} cujus pars %
\textit{(aaaaa-a)}~\textit{OPQR} %
\textit{(bbbbb-b)}~\textit{MPQR}%
~\textit{L}}}%
sit tubus\indexS{tubus} rectilineus %
\edtext{clausus in \textit{M}}{\lemma{}\Bfootnote{clausus in \textit{M}~\textit{erg.~L}}}
aequalis ubique amplitudinis, similiter %
\edtext{et pars \textit{STV} sit tubus\indexS{tubus}}{\lemma{}\Bfootnote{et \lbrack...\rbrack\ tubus~\textit{erg.~L}}}
priori parallelus et capacitate ubique aequalis, cui in \textit{V} insistat pila cava\indexS{pila cava} \textit{VXZ},
\edtext{cum tubulo\indexS{tubus}}{\lemma{}\Bfootnote{cum \ \textbar \ exiguo \textit{ gestr.}\ \textbar\ tubulo \textit{ L}}} %
\edtext{aperto}{\lemma{}%
\Bfootnote{aperto~\textit{erg.~L}}} %
in
\edtext{\textit{Z}, intelligantur autem \textit{QR} vel \textit{ST} aequales et \textit{PQ} aequalis sit}{%
\lemma{\textit{Z},}%
\Bfootnote{%
\textit{(1)}~\bars{ arcto quidem sed satis longo, ut sensu notabili \textit{erg. u. gestr.}}
intelligantur \textit{OP}, \textit{PQ}, \textit{QR}, \textit{ST}, \textit{TV} aequales, et quaevis %
\bars{pars \textit{streicht Hrsg.}}
trium circiter pedum, ita ut %
\textit{(2)}~intelligantur autem %
\textit{(a)}~aequales \textit{PR}, et \textit{SV}, item \textit{QR} et \textit{ST} quae duae portiones aequentur %
\textit{(aa)}~(quaelibet) %
\textit{(bb)}~quae et %
\textit{(b)}~\textit{QR} vel \lbrack...\rbrack\ aequalis sit %
~\textit{L}}}
altitudini hydrargyri\indexS{hydrargyrum} in tubo %
\edtext{Torricelliano\indexS{tubus Torricellianus}. Dum ergo immittatur tantum Hydrargyri\indexS{hydrargyrum}}{\lemma{Torricelliano}\Bfootnote{%
\textit{(1)}~Immisso %
\textit{(a)}~deinde %
\textit{(b)}~jam Hydrargyro %
\textit{(2)}~. Dum %
\textit{(a)}~ja %
\textit{(b)}~venditum %
\textit{(c)}~ergo immittatur tantum Hydrargyri~\textit{L}}}
per %
\edtext{\textit{Z} quantum capere potest \textit{PQRST},}{\lemma{}\Bfootnote{quantum capere potest \textit{PQRST}~\textit{erg.~L}}}
ita inclinetur tubus\indexS{tubus}, %
ut ab Hydrargyro\indexS{hydrargyrum} impleatur usque %
\edtext{in \textit{M} quo facto}{\lemma{in}\Bfootnote{%
\textit{(1)}~\textit{O} %
\textit{(2)}~\textit{M}. %
\textit{(a)}~Igne %
\textit{(aa)}~deinde %
\textit{(bb)}~denique %
\textit{(b)}~Deinde igne lampadis claudatur tubus vel hermetice sigillatur in
\textit{(aa)}~\textit{O} %
\textit{(bb)}~\textit{M} %
\textit{(3)}~quo facto
~\textit{L}}}
\edtext{tubus\indexS{tubus}}{\lemma{}\Bfootnote{tubus~\textit{erg.~L}}}
erigatur %
\edtext{ut \textit{MR}}{\lemma{ut}\Bfootnote{%
\textit{(1)}~\textit{OR} %
\textit{(2)}~\textit{MR}~\textit{L}}}
sit %
\edtext{perpendicularis ad horizontem}{\lemma{perpendicularis}\Bfootnote{%
\textit{(1)}~, et extre %
\textit{(2)}~ad horizontem~\textit{L}}}
et %
\edtext{extremitas \textit{M}}{\lemma{extremitas}\Bfootnote{%
\textit{(1)}~\textit{O} %
\textit{(2)}~\textit{M}~\textit{L}}}
superior; decidetque %
Argentum vivum usque in \textit{PQRST},
\edtext{et pars tubi\indexS{tubus} \textit{MP}}{\lemma{et}\Bfootnote{%
\textit{(1)}~in \textit{MP} %
\textit{(2)}~pars tubi \textit{MP}~\textit{L}}}
manebit aere vacua. %
\edtext{Claudatur nunc exiguus}{\lemma{Claudatur}\Bfootnote{%
\textit{(1)}~jam %
\textit{(2)}~nunc \bars{igne lampadis \textit{ gestr.}} exiguus~\textit{L}}}
quoque tubulus\indexS{tubus} in \textit{Z} %
\href{\Scan{LH_37_04_169v}}%
{\lbrack169\:v\textsuperscript{o}\rbrack}\ %
et paratum erit \textls{Thermometrum\indexS{thermometrum} clausum} quod %
aeris ambientis pondere\indexS{pondus aeris} non potest %
\edtext{mutari. Proinde convenit autem hydrargyrum\indexS{hydrargyrum}}{\lemma{mutari.}\Bfootnote{%
\textit{(1)}~Utile autem est %
\textit{(2)}~Proinde convenit autem hydrargyrum~\textit{L}}}
prius fuisse ab aere\indexS{aer} purgatum 
quod fit si aliquamdiu steterit 
\edlabel{LH_37_04_169-170_00a}
in machina\indexS{machina}
\edtext{ab insigni viro}{\lemma{}\Bfootnote{ab insigni viro~\textit{erg.~L}}}
Ottone Gerickio\indexN{\textls{Guericke} (Gerickius, Gericke), Otto von 1602\textendash1686} 
primum%
\edtext{Magdeburgi pulchre excogitata\edlabel{LH_37_04_169-170_00b}}{\lemma{Magdeburgi}
\Bfootnote{%
\textit{(1)}~inventa %
\bars{
\textit{(2)}~ingeniosissime %
\textit{(3)}~felici %
\textit{(4)}~pulchre excogitata~\textit{erg.}}~\textit{L}}}%
\edtext{}{%
\xxref{LH_37_04_169-170_00a}{LH_37_04_169-170_00b}%
\lemma{in \lbrack...\rbrack\ excogitata}%
\Cfootnote{\textsc{O.~von Guericke}, \textit{Experimenta nova}, lib.~III, cap. IV, Amsterdam 1672, S.\,75\textendash77.\cite{00055}}}%
\edlabel{LH_37_04_169-170_01a}
et a celeberrimo 
\edlabel{LH_37_04_169-170_02a}
Boylio\indexN{\textls{Boyle} (Boylius, Boyl), Robert 1627\textendash1691} egregie exculta, ex qua aer\indexS{aer} suctu est exhaustus. Utile quoque est %
\edlabel{LH_37_04_169-170_02b}%
\edtext{}{%
\xxref{LH_37_04_169-170_02a}{LH_37_04_169-170_02b}%
\lemma{Boylio}%
\Bfootnote{%
\textit{(1)}~promota (unde aer suctu exhauritur) et %
\textit{(2)}~,~egregie \lbrack...\rbrack\ est~\textit{L}}}%
\edlabel{LH_37_04_169-170_01b}%
\edtext{}{%
\xxref{LH_37_04_169-170_01a}{LH_37_04_169-170_01b}%
\lemma{et \lbrack...\rbrack\ exculta}%
\Cfootnote{\textsc{R. Boyle}, \textit{New experiments physico-mechanicall}, Oxford 1660, S.\,1\textendash19.\cite{00015}}}%
\edtext{thermometrum\indexS{thermometrum} praeparari tempore quo}{\lemma{thermometrum}\Bfootnote{%
\textit{(1)}~fabricari %
\textit{(2)}~praeparari %
\textit{(a)}~vel saltem %
\textit{(b)}~tempore quo~\textit{L}}}
medius est circiter tam caloris\indexS{calor} in aere\indexS{aer}, quam %
ponderis in atmosphaera\indexS{atmosphaera} status. Ponamus jam caloris\indexS{calor} %
\edtext{in aere\indexS{aer} \textit{TVXZ}}{\lemma{}\Bfootnote{%
in aere \textit{TVXZ}~\textit{erg.~L}}}
\edtext{gradum\indexS{gradus} adeoque et dilatandi se vim mutari et vel ea aucta}{\lemma{gradum}\Bfootnote{%
\textit{(1)}~mutari et vel eo aucto %
\textit{(2)}~adeoque \lbrack...\rbrack\ aucta~\textit{L}}}
Hydrargyrum\indexS{hydrargyrum} ex \textit{P} %
ascendere in \textit{F}, vel minuta eadem ex \textit{P} descendere %
in \textit{\greek{f}}.
\edtext{Quaeritur quis}{\lemma{Quaeritur}\Bfootnote{%
\textit{(1)}~quomodo gradus %
\textit{(2)}~quis~\textit{L}}}
in casu hydrargyri\indexS{hydrargyrum} positi %
\edtext{in \textit{F} vel \textit{\greek{f}}}{\lemma{in}\Bfootnote{%
\textit{(1)}~\textit{P} ve %
\textit{(2)}~\textit{F} vel \textit{\greek{f}}~\textit{L}}}
sit caloris gradus\indexS{gradus caloris} %
\edtext{posita ratione rectae \textit{PF} vel \textit{P\greek{f}} ad \textit{PQ} datam,}{\lemma{}\Bfootnote{posita
ratione \bars{\textit{(1)} rectarum \textit{(2)} rectae} \textit{PF} \lbrack...\rbrack\ datam~\textit{erg.~L}}}
seu quis numerus %
\edtext{divisionis ascribendus sit}{\lemma{divisionis}\Bfootnote{%
\textit{(1)}~esset %
\textit{(2)}~ascribendus sit~\textit{L}}}
in \textit{F} vel \textit{\greek{f}}.
Hoc ita vestigabimus:
\pend
\vspace{0.75em}
\pstart\noindent
\lbrack \textit{Nachfolgend kleingedruckter Text in L gestrichen:}\rbrack%
\pend
\vspace{0.5em}
\pstart
\normalsize
\centering{%
\lbrack \textit{Erste Fassung}\rbrack}
\pend 
\pstart
\footnotesize
Capacitas %
totius instrumenti\indexS{instrumentum} sit \textit{A} et %
\edtext{capacitas seu spatium}{\lemma{}\Bfootnote{capacitas seu spatium \textit{ erg.} \textit{ L}}} %
Hydrargyri\indexS{spatium hydrargyri} 
\edtext{infusi (seu portionis \textit{PQRST}) sit \textit{G};}{%
\lemma{infusi}%
\Bfootnote{%
\textit{(1)}~\textit{G}, et
\textit{(2)}~(seu \lbrack...\rbrack\ sit \textit{G}~\textit{L}}}
et spatium seu %
\edtext{volumen Hydrargyri\indexS{volumen hydrargyri}}{\lemma{volumen}\Bfootnote{%
\textit{ (1) } Mercurii\indexS{mercurius} %
\textit{ (2) } Hydrargyri suspensi~\textit{L}}}
suspensi \textit{PQ} sit \textit{H}. %
\edtext{Ergo volumen}{\lemma{}\Bfootnote{Ergo \ \textbar \ et spa \textit{erg. u. gestr.}\ \textbar\ volumen \textit{ L}}}
aeris\indexS{aer} 
\edtext{\textit{TVXZ}}{\lemma{}\Bfootnote{\textit{TVXZ} \textit{ erg.} \textit{ L}}}
habebitur, si a %
\edtext{spatio \textit{QRSTVXZ}}{\lemma{spatio}\Bfootnote{%
\textit{ (1) } toto \textit{A} %
\textit{ (2) } \textit{QRSTVXZ}~\textit{L}}}
detrahatur %
\edtext{\textit{G}}{\lemma{}\Bfootnote{\textit{G} \textit{ erg.} \textit{ L}}}
spatium hydrargyri\indexS{spatium hydrargyri} %
\edtext{totius \textit{QRST}.}{\lemma{totius}\Bfootnote{%
\textit{(1)}~\textit{G} %
\textit{(2)}~\textit{QRST}.
~\textit{L}}}
Sed %
spatium \textit{QRSTVXZ} habetur si a toto spatio \textit{A} detrahatur 
\edtext{spatium \textit{OP},}{%
\lemma{spatium \textit{OP}}%
\applabel{LH_37_04_169-170_21ab}%
\Cfootnote{Der Buchstabe \textit{O} bezeichnet hier wie an weiteren Stellen des gestrichenen Textes den Punkt \textit{M} in Fig.~2. Die Bezeichnung \textit{O} geht auf eine frühe Fassung der Zeichnung zurück, wie sie noch auf \href{\Scan{LH_037_05_108r}}{Bl.~108\:r\textsuperscript{o}} von N.~\ref{60326} zu erkennen ist.}}
id est \lbrack\textit{Text bricht ab.}\rbrack%

\pend
\vspace{0.5em}
\pstart
\normalsize
\centering{%
\lbrack \textit{Zweite Fassung}\rbrack}
\pend 
\pstart
\footnotesize
Capacitas ipsius instrumenti\indexS{instrumentum} integri %
seu spatii \textit{OPQRSTVXZ} sit \textit{A}; spatium %
ipsius hydrargyri\indexS{spatium hydrargyri}, ut \textit{PQRST} %
\edtext{vel \textit{FGRSC},}{\lemma{}\Bfootnote{vel \textit{FGRSC} \textit{ erg.} \textit{ L}}}
sit \textit{B}. %
Haec duo %
\edtext{\textit{A} et \textit{B}}{\lemma{}\Bfootnote{\textit{A} et \textit{B}~\textit{erg.~L}}}
semper manent eadem. %
\pend
\vspace{1.25em}
\pstart
\normalsize
\centering{%
\lbrack \textit{Erste Fassung}\rbrack}
\pend 
\pstart
\footnotesize
Constat autem spatium %
\edtext{hydrargyri\indexS{spatium hydrargyri} \textit{PQRST} ex spatio}{\lemma{hydrargyri}\Bfootnote{%
\textit{ (1) } ex spatio \textit{RS} semper manente %
\textit{ (2) } \textit{PQRST} ex spatio~\textit{L}}} %
\textit{RS} quod sit $\varepsilon$ et spatio \textit{ST} duplicato, (nempe %
\textit{QR} et \textit{ST}) quod sit bis \textit{D}; et spatio \textit{PQ} %
\edtext{Hydrargyri\indexS{spatium hydrargyri} ab aere\indexS{aer} ordinario sustentati}{\lemma{}\Bfootnote{Hydrargyri \lbrack...\rbrack\ sustentati \textit{ erg.} \textit{ L}}} %
quod %
est datum, data quippe altitudine \textit{PQ} eadem %
quae est in Tubo Torricelliano\indexS{tubus Torricellianus}. Habemus %
ergo \textit{B} aequal. $\varepsilon$ %
\edtext{plus}{\lemma{}\Bfootnote{plus \textit{ erg.} \textit{ L}}} %
bis \textit{D} %
\edtext{plus}{\lemma{}\Bfootnote{plus \textit{ erg.} \textit{ L}}} %
\textit{H}. Similiter constat \lbrack\textit{Text bricht ab.}\rbrack
\pend
\vspace{0.25em}
\pstart
\normalsize
\centering{%
\lbrack \textit{Zweite Fassung}\rbrack}
\pend 
\pstart
\footnotesize
\edtext{Constat autem}{\lemma{}\Bfootnote{Constat autem \textit{ erg.} \textit{ L}}} %
spatium hydrargyri\indexS{spatium hydrargyri} \textit{FGRSC} ex spatio \textit{RS} seu $\varepsilon$, %
spatio \textit{SC} duplicato (nempe \textit{GR} et \textit{SC}) quod sit %
bis \textit{L}, et \textit{FG} quod sit \textit{M}, spatium scilicet Hydrargyri\indexS{spatium hydrargyri} %
ab aere\indexS{aer} nunc alterato sustentati; et habemus %
iterum \textit{B} aequal. $\varepsilon$ %
\edtext{plus}{\lemma{}\Bfootnote{plus \textit{ erg.} \textit{ L}}} %
bis \textit{L}
\edtext{plus}{\lemma{}\Bfootnote{plus \textit{ erg.} \textit{ L}}} %
\textit{M}. Ergo %
\edtext{fit \textit{B} minus $\varepsilon$ aequal.}{%
\lemma{fit}\Bfootnote{%
\textit{(1)}~bis \textit{D} plus \textit{M} %
\textit{(2)}~\textit{B} minus $\varepsilon$ 
\textit{(a)}~seu %
\textit{(b)}~aequal.%
~\textit{L}}}
bis %
\edtext{\textit{L} plus \textit{M}.}{\lemma{\textit{L}}%
\Bfootnote{%
\textit{(1)}~+~\textit{M} %
\textit{(2)}~plus \textit{M}~\textit{L}}}
Porro %
spatium aeris\indexS{spatium aeris} ante alterationem, nempe \textit{TVXZ} %
sit \textit{N}%
, cui si addatur in %
\edtext{\textit{TC} in casu}{%
\lemma{\textit{TC}}%
\Bfootnote{%
\textit{(1)}~vel auferatur %
\textit{(2)}~in casu~\textit{L}}}
aucti caloris\indexS{calor}, (vel auferatur \textit{TK} in casu %
minuti) 
\edlabel{LH_37_04_169-170_03a}%
fit \lbrack\textit{Text bricht ab.}\rbrack%
\pend
\pstart
\footnotesize
\edtext{Datur et \textit{PQ} spatium Hydrargyri\indexS{spatium hydrargyri} ab aere\indexS{aer} ante alterationem sustentati, cujus altitudo}{%
\xxref{LH_37_04_169-170_03a}{LH_37_04_169-170_03b}%
\lemma{fit}%
\Bfootnote{%
\textit{(1)}~Spatium \textit{PQ} %
\textit{(2)}~Datur et \textit{PQ} spatium %
\textit{(a)}~ex quo \textit{M} %
\textit{(b)}~Hydrargyri ab aere %
\textit{(aa)}~ordinario %
\textit{(bb)}~ante alterationem sustentati, %
\textit{(aaa)}~quod idem est cum alti %
\textit{(bbb)}~cujus altitudo %
~\textit{L}}}%
\edlabel{LH_37_04_169-170_03b} %
est eadem cum altitudine ejusdem in tubo Torricelliano\indexS{tubus Torricellianus} %
suspensi, quod vocetur \textit{D}; %
\edlabel{LH_37_04_169-170_04a} 
itemque datur
\href{\Scan{LH_37_04_170r}}%
{\lbrack170\:r\textsuperscript{o}\rbrack}\
\textit{TVXZ} spatium %
aeris\indexS{aer} ante alterationem, %
quod vocetur %
$\varepsilon$ %
jam \textit{CT} vel \textit{RQ} vocetur \textit{L}.
\edlabel{LH_37_04_169-170_04b}
\edtext{}{%
\xxref{LH_37_04_169-170_04a}{LH_37_04_169-170_04b}%
\lemma{datur}%
\Bfootnote{%
\textit{(1)}~$\varepsilon$ spatium %
\textit{(2)}~\textit{TVXZ} spatium, \lbrack...\rbrack\ vocetur $\varepsilon$, jam %
\textit{(a)}~\textit{RS} vocetur \textit{H} et \textit{C}\textlangle\textit{Q}\textrangle\ vel %
\textit{(b)}~\textit{CT} vel \textit{RQ} vocetur \textit{L}.~\textit{L}}} %
\pend
\pstart
\footnotesize
\edtext{Habemus }{\lemma{}\Bfootnote{Habemus \textit{ erg.~L}}}%
\textit{B} aequal. \textit{D} plus \textit{RS} plus bis %
\edtext{\textit{L}, seu bis \textit{L} aequal. \textit{B} minus \textit{D} minus \textit{RS}.}{\lemma{\textit{L},}%
\Bfootnote{%
\textit{(1)}~et \textit{A} %
\textit{(2)}~seu \lbrack...\rbrack\ \textit{RS}.%
~\textit{L}}}
\edtext{Porro \textit{A} aequal. \textit{OR}}{\lemma{Porro}%
\Bfootnote{%
\textit{(1)}~\textit{OR} vocemus \textit{M}, fit %
\textit{(2)}~\textit{A} aequal. %
\textit{(a)}~\textit{M} %
\textit{(b)}~\textit{OR}%
~\textit{L}}}
plus \textit{RS} %
plus \textit{L} plus%
\edtext{$\varepsilon$ seu }{\lemma{$\varepsilon$} %
\Bfootnote{%
\textit{(1)}~similiter in casu aeris %
\textit{(2)}~seu%
~\textit{L}}}%
bis \textit{L} aequal. bis \textit{A} minus bis \textit{OR} minus bis \textit{RS} %
minus bis %
\edtext{$\varepsilon$ seu \textit{B} }{\lemma{$\varepsilon$}%
\Bfootnote{%
\textit{(1)}~. Ergo fit bis %
\textit{(2)}~seu \textit{B}%
~\textit{L}}}%
minus \textit{D} minus %
\edtext{\textit{RS}. Ergo fit \textit{D} plus bis $\varepsilon$ aequal. \textit{B} minus bis \textit{A} plus \textit{ORS}. Similiter igitur post alterationem aeris\indexS{aer} translato }{\lemma{\textit{RS}}%
\Bfootnote{%
\textit{(1)}~Ergo fit \textit{D} aequal. \textit{B} plus \textit{OR} \bars{\textit{ORS}~\textit{erg. u. gestr.}} plus \textit{RS} minus bis \textit{A} minus bis $\varepsilon$. %
\textit{(a)}~Spatium igitur aeris post %
\textit{(b)}~igitur
\textit{(2)}~Ergo fit %
\textit{(a)}~bis $\varepsilon$ minus \textit{D} aequal. %
\textit{(b)}~\textit{D} plus \lbrack...\rbrack\ aeris %
\textit{(aa)}~altit %
\textit{(bb)}~spatium %
\textit{(cc)}~translato~\textit{L}}}%
Hydrargyro\indexS{hydrargyrum} ex \textit{PQRST} in %
\textit{FGRSC}, et altitudinem %
\edtext{Hydrargyri\indexS{hydrargyrum} \textit{FG}}{\lemma{Hydrargyri}%
\Bfootnote{%
\textit{(1)}~\textit{FQ} %
\textit{(2)}~\textit{FG}%
~\textit{L}}} %
quam nunc aer\indexS{aer} sustinet vocando \textit{M} 
\edtext{et aeris spatium\indexS{spatium aeris} \textit{CVXZ}}{\lemma{et}%
\Bfootnote{%
\textit{(1)}~spatium quod \bars{nunc occupat~\textit{streicht Hrsg.}} %
\textit{(2)}~aeris spatium \textit{CVXZ}%
~\textit{L}}} %
vocando \textit{N}, fit %
\textit{M} plus bis \textit{N} %
\edtext{aequal. \textit{B} minus bis \textit{A}}{\lemma{aequal.}%
\Bfootnote{%
\textit{(1)}~\textit{D} plus bis $\varepsilon$ %
\textit{(2)}~\textit{B} minus bis \textit{A}
~\textit{L}}} %
plus \textit{ORS}, itaque \textit{M} plus bis \textit{N} aequal. \textit{D} plus bis $\varepsilon$
\edlabel{LH_37_04_169-170_05a}
jam \lbrack\textit{Text bricht ab.}\rbrack%
\pend%
\vspace{0.5em} 
\normalsize
\pstart 
\edtext{Hydrargyro\indexS{hydrargyrum}}{%
\xxref{LH_37_04_169-170_05a}{LH_37_04_169-170_05b}%
\lemma{jam}%
\Bfootnote{%
\textit{(1)}~Si %
\textit{(2)}~Hydrargyro%
~\textit{L}}} %
 \edlabel{LH_37_04_169-170_05b}
ex %
\edtext{\textit{PQRST} (ubi pars cujus \textit{PQ} super aequilibrium \textit{QRST} suspensa fuerat) 
\edlabel{LH_37_04_169-170_07a}
per auctum calorem\indexS{calor}
\edlabel{LH_37_04_169-170_07b}
aeris\indexS{aer} (ex \textit{ZXVT} se expandentis in \textit{ZXVC}) translato in \textit{FGRSC} (ita ut pars Hydrargyri\indexS{hydrargyrum} \textit{FG} jam suspensa sit super aequilibrium \textit{GRSC}) patet esse \textit{TC} aequal.\,\textit{QG} et rursus \textit{TC} aequal.\,\textit{PF}.}%
{\lemma{\textit{PQRST}}%
\Bfootnote{%
\textit{(1)}~translato in \textit{FGRSC} %
\textit{(2)}~vi caloris aeris in \textit{ZVT} se expandentis in \textit{ZVC} translato \bars{ubi~\textit{erg. u. gestr.}} patet esse %
\textit{(a)}~\textit{CT} %
\textit{(aa)}~et %
\textit{(bb)}~aeq %
\textit{(b)}~\textit{TC} aequal. \textit{QG}, rursus \textit{TC} aequal. \textit{PF}. %
\textit{(3)}~(ubi pars cujus \textit{PQ} super aequilibrium \textit{QRST} suspensa fuerat) %
\textit{(a)}~\textlangle vi\textrangle\ aucti caloris (aeris ex \textit{ZXV} \bars{se~\textit{streicht Hrsg.}} expandentis \bars{in~\textit{streicht Hrsg.}} \textit{ZXVT}) %
\textit{(b)}~per vim aucti caloris in aere %
\textit{(c)}~per auctum calorem aeris \lbrack...\rbrack\ \textit{PF}.
~\textit{L}}} %
\edtext{}{%
\xxref{LH_37_04_169-170_07a}{LH_37_04_169-170_07b}%
\lemma{}%
\Afootnote{%
\textit{Am Rand, gestr.}: \footnotesize{per calo}%
}}%
Ergo \textit{GF} minus \textit{QP} aequal. bis %
\edtext{\textit{TC}\lbrack,\rbrack\ id est differentia}{\lemma{\textit{TC}}%
\Bfootnote{%
\textit{(1)}~est autem \textit{GF} %
\textit{(2)}~.~\textit{GF} minus \textit{QP} differentia %
\textit{(3)}~id est dupla %
\textit{(4)}~id est differentia%
~\textit{L}}}
inter %
\edtext{duo loca seu volumina Hydrargyri\indexS{hydrargyrum}}{\lemma{duo}%
\Bfootnote{%
\textit{(1)}~spatia %
\textit{(2)}~volumina %
\textit{(3)}~altitudines Mercu %
\textit{(4)}~loca seu volumina Hydrargyri%
~\textit{L}}}
sustentati in diversis aeris\indexS{aer} secundum %
calorem\indexS{calor} statibus seu inter \textit{GF} et \textit{QP}, est %
\edtext{aequalis ipsi \textit{C}\lbrack\textit{K}\rbrack, duplae}{%
\lemma{aequalis}%
\Bfootnote{%
\textit{(1)}~duplae %
\textit{(2)}~ipsi \bars{\textit{CT~L ändert Hrsg.}}, duplae%
~\textit{L}}}
differentiae inter volumina aeris\indexS{volumen aeris} ejusdem %
\edtext{calore\indexS{calor} diversi,}{\lemma{calore}%
\Bfootnote{%
\textit{(1)}~differentis %
\textit{(2)}~diversi \bars{ \textit{CT}~\textit{gestr.}}%
~\textit{L}}}
idemque concluditur si minuto %
calore\indexS{calor} Hydrargyrum\indexS{hydrargyrum} ex \textit{PQRT} %
\edtext{in \textit{\greek{f}HRSK} translatum ponamus. Nam}{\lemma{}%
\Bfootnote{%
in \bars{\textit{\greek{f}HRSK}~\textit{erg., streicht Hrsg.}} \textit{\greek{f}HRSK} %
\textit{(1)}~ad aequilibrium \textit{DRSV} magis accessisse ponamus, (a quo tamen semper longe abest, %
\textit{(a)}~itaque %
\textit{(b)}~et proinde sufficit \textit{V} initium pilae esse ex adverso ipsius \textit{H} puncti medii inter \textit{P} et \textit{Q}) 
\textit{(2)}~translatum ponamus. Nam%
~\textit{L}}} %
$QP-H\text{\textit{\greek{f}}}$ aequal. bis %
\edtext{\textit{C}\lbrack\textit{T}\rbrack}{%
\lemma{}%
\Bfootnote{%
 \textit{CK}
\textit{L ändert Hrsg.}}}
et vero minuto %
\edtext{calore\indexS{calor} Mercurius\indexS{mercurius} accedit}{\lemma{calore}%
\Bfootnote{%
\textit{(1)}~sem %
\textit{(2)}~ab Hydrargyro %
\textit{(3)}~Hydrargyrum acceditur %
\textit{(4)}~Mercurius accedit%
~\textit{L}}}
ad aequilibrium \textit{DRSV}, quanquam %
\edtext{id nunquam attingat, et multum semper inde absit.}{\lemma{quanquam}%
\Bfootnote{%
\textit{(1)}~ longe semper ab eo %
\textit{(2)}~id \lbrack...\rbrack\, inde %
\textit{(a)}~desit %
\textit{(b)}~absit%
~\textit{L}}}
(Itaque cum \textit{D} sit %
punctum medium inter \textit{P} et \textit{Q}, quorum distantia determinata %
est, potest pila\indexS{pila} in \textit{V} incipere ex %
adverso ipsius \textit{D}.)
\pend
\pstart
\edlabel{LH_37_04_169-170_18a}%
Jam %
\edtext{generaliter sine discrimine caloris\indexS{calor} aucti vel minuti}{\lemma{}\Bfootnote{generaliter \lbrack...\rbrack\ minuti~\textit{erg.~L}}}
volumen aeris\indexS{volumen aeris} ante alterationem %
\edtext{nempe \textit{ZXVT} numero expressum}{\lemma{}\Bfootnote{nempe \textit{ZXVT} %
\bars{%
\textit{(1)} numeri ratione %
\textit{(2)} numero} %
expressum%
~\textit{erg.~L}}}
vocemus \textit{A}, post %
\edtext{alterationem nempe \textit{ZXVC} vel \textit{ZXVK} vocemus \textit{B}; et altitudinem seu volumen ante alterationem, nempe \textit{PQ} vocemus \textit{L}, post eam nempe \textit{FG} vel \textit{\greek{f}H} vocemus \textit{M}. Differentiam autem inter \textit{A} et \textit{B} %
(quae est \textit{TC} vel \textit{TK}) vocemus \textit{N}, %
erit differentia inter \textit{L} et \textit{M} aequalis duplae \textit{N}. %
\edlabel{LH_37_04_169-170_18b}%
Dantur autem \textit{L} et \textit{M}, et \textit{N}, ex inspectione Tubi\indexS{tubus} aequabiliter crassi, cujus sunt partes, datur et \textit{A}, seu capacitas ipsius spatii \textit{ZXVT} %
quia %
capacitas pilae\indexS{pila}, seu partis instrumenti\indexS{instrumentum} inaequaliter crassae, sive ejus ratio ad tubum\indexS{tubus} aequaliter crassum vel partes ejus altitudine datas, potuit antequam clauderetur tubus\indexS{tubus}, aqua immissa explorari. Datur ergo et \textit{B} seu $A+N$ vel $A-N$.%
}{%
\lemma{alterationem}%
\Bfootnote{%
\textit{(1)}~\textit{B} \bars{et~\textit{streicht Hrsg.}} volumen %
\textit{(a)}~Mercu %
\textit{(b)}~\bars{Hydrargyri sustentati~\textit{streicht Hrsg.}} ante alterationem vocemus %
\textit{(aa)}~\textit{D}, post $\varepsilon$ %
\textit{(bb)}~\textit{L} post eam \textit{M}, erit differentia inter %
\textit{(aaa)}~\textit{D} et $\varepsilon$ %
\textit{(bbb)}~\textit{L} et \textit{M} aequalis duplae differentiae inter \textit{A} et \textit{B}. Sed rursus %
\textit{(aaaa)}~\textit{TVXZ} (id est \textit{A}) plus \textit{TC} plus %
\textit{(bbbb)}~\textit{A} seu \textit{TVXZ}, plus \textit{CT}, plus %
\textit{(2)}~nempe \textit{ZXVC} vel \textit{ZXVK} vocemus %
\textit{(a)}~\textit{A} %
\textit{(b)}~\textit{B}; \lbrack...\rbrack\ vocemus \textit{M}
\textit{(aa)}~, erit differentia inter \textit{L} et \bars{\textit{M} aequalis~\textit{streicht Hrsg.}} duplae differentiae inter \textit{A} et \textit{B}. %
\textit{(bb)}~.~Datur jam \textit{L} et \textit{M} ex differentia %
\textit{(cc)}~.~Differentiam autem inter \textit{A} et \textit{B} %
\textit{(aaa)}~nempe \textit{TC} vel \textit{CK} %
\textit{(bbb)}~(quae est \textit{TC} vel \textit{TK}) vocemus \textit{N}, %
\textit{(aaaa)}~fiet %
\textit{(bbbb)}~erit differentia \lbrack...\rbrack\ spatii \textit{ZXVT} %
\textit{(aaaaa)}~cujus %
\textit{(bbbbb)}~quia %
\textit{(aaaaa-a)}~magnitudo %
\textit{(bbbbb-b)}~capacitas pilae, \lbrack...\rbrack\ vel $A-N$.%
~\textit{L}}} %
\pend 
\vspace{0.75em}
\pstart
\noindent
\lbrack \textit{Nachfolgend kleingedruckter Text in L gestrichen:}\rbrack%
\pend
\vspace{0.25em}
\pstart 
\footnotesize
Porro volumen totius instrumenti\indexS{instrumentum} \textit{ORSVXZ} vocemus \textit{L} %
et volumen totius Hydrargyri\indexS{hydrargyrum} vocemus \textit{M}, %
\edtext{et \textit{OP} sit \textit{N.}}{\lemma{}\Bfootnote{et \textit{OP} sit \textit{N.}~\textit{erg.~L}}}
Patet fore %
\textit{A} (hoc est \textit{TVXZ}), plus \textit{CT}, plus \textit{M} (seu \textit{CSRF}) plus %
\textit{OF} (id est %
\edtext{plus \textit{OP} minus}%
{\lemma{plus}%
\Bfootnote{%
\textit{(1)}~\textit{PO} min %
\textit{(2)}~\textit{OP} minus%
~\textit{L}}}
\textit{PF} seu plus \textit{N} minus \textit{CT}) %
aequal. \textit{L}.
\edlabel{LH_37_04_169-170_13a}%
\normalsize
\href{\Scan{LH_37_04_170v}}%
{\lbrack170\:v\textsuperscript{o}\rbrack}\
\pend
\pstart
\footnotesize
\edtext{Datur vero ratio voluminis \textit{TVXZ}, %
\edlabel{LH_37_04_169-170_13b}%
}{%
\xxref{LH_37_04_169-170_13a}{LH_37_04_169-170_13b}%
\lemma{\textit{L.}}%
\Bfootnote{%
\textit{(1)}~Detur jam ratio %
\textit{(2)}~Datur %
\textit{(a)}~jam %
\textit{(b)}~vero ratio voluminis %
\textit{(aa)}~esse %
\textit{(bb)}~\textit{TVXZ}%
~\textit{L}}}
seu ipsius \textit{L} %
ad \textit{PQ} 
\edtext{sive \textit{A}}{\lemma{}\Bfootnote{sive \textit{A}~\textit{erg.~L}}}
volumen hydrargyri\indexS{hydrargyrum} ab aere\indexS{aer} nondum alterato %
\edtext{sustentati. Nam divisa}{\lemma{sustentati,}%
\Bfootnote{%
\textit{(1)}~seu ad %
\textit{(2)}~.~Nam divisa%
~\textit{L}}}
altitudine \textit{PQ} in %
1000 %
\edtext{(si placet)}{\lemma{}\Bfootnote{(si placet)~\textit{erg.~L}}}
partes aequales %
\edtext{ita ut \textit{A} sit 1000,}{\lemma{}\Bfootnote{ita ut \textit{A} sit 1000~\textit{erg.~L}}}
numerus %
\edtext{partium quibus aequivalet \textit{TVXZ} si erit}%
{\lemma{partium}%
\Bfootnote{%
\textit{(1)}~quas continet %
\textit{(2)}~quibus %
\textit{(a)}~aequivalent %
\textit{(b)}~aequivalet %
\textit{(aa)}~sit %
\textit{(bb)}~si erit%
~\textit{L}}}
ipse numerus \textit{L} %
\edtext{isque utique facile sciri potest implendo %
pilam\indexS{pila} \textit{VXZ} aqua, et %
\edtext{eandem aquam}{\lemma{}%
\Bfootnote{%
eandem \bars{deinde~\textit{gestr.}} aquam%
~\textit{ L}}}
in tubum\indexS{tubus} transfundendo, atque %
observando quot partes tubi\indexS{tubus}
\edtext{aqua}{\lemma{}%
\Bfootnote{aqua~\textit{gestr. L, wieder gültig gemacht Hrsg.}}} 
occupet, quibus deinde addentur partes %
\textit{VT}, ut habeatur 
\edlabel{LH_37_04_169-170_09a}
\textit{L}.
}{\lemma{}%
\Bfootnote{isque \lbrack...\rbrack\ habeatur \textit{L}~\textit{erg.~L}}}
\edtext{Porro}{%
\xxref{LH_37_04_169-170_09a}{LH_37_04_169-170_09b}%
\lemma{\textit{L}.}%
\Bfootnote{%
\textit{(1)}~,~cui si addetur et %
\textit{(2)}~.~Porro%
~\textit{L}}}
\edlabel{LH_37_04_169-170_09b}%
diviso in similes partes toto %
\edtext{instrumento\indexS{instrumentum} \textit{MRSV},}%
{\lemma{instrumento}%
\Bfootnote{%
\textit{(1)}~\textit{VSRO} %
\textit{(2)}~\textit{MRSV}%
~\textit{L}}}
dabitur \textit{TC} %
\edtext{vel \textit{PF},}{\lemma{}%
\Bfootnote{vel \textit{PF}~\textit{erg.~L}}}
ex inspectione %
voceturque \textit{N}; (eodemque modo \textit{P\greek{f}}) %
\edtext{%
\edlabel{LH_37_04_169-170_10a}%
itaque et datur \textit{M} aequale \textit{L} plus (vel minus) \textit{N}.%
}%
{%
\xxref{LH_37_04_169-170_10a}{LH_37_04_169-170_10b}%
\lemma{itaque}%
\Bfootnote{%
\textit{(1)}~\textit{M} %
\textit{(a)}~est %
\textit{(b)}~aequal. \textit{L} plus (vel minus) \textit{N}, datur. %
\textit{(2)}~et datur \textit{M} aequal. \textit{L} plus (vel minus) \textit{N} %
~\textit{L}}}%
\edlabel{LH_37_04_169-170_10b}%
\edtext{Similiter ex inspectione datur \textit{B} seu numerus partium in \textit{GF} et \textit{B} seu numerus partium earundem quas continet \textit{FG} est aequal. \textit{A} plus (vel minus) bis \textit{N}.}
{\lemma{}
\Bfootnote{Similiter \bracks{...} aequal. \textit{A} \bars{\textit{(1)}~(seu \textit{(PO}) \textit{(2)}~plus (vel minus) bis \textit{N}} \textit{erg.~L}}}
\edlabel{LH_37_04_169-170_26a}%
\edtext{Porro densitates aeris\indexS{densitas aeris} \bracks{\textit{Text bricht ab}}}{%
\xxref{LH_37_04_169-170_26a}{LH_37_04_169-170_26b}%
\lemma{Porro}%
\Bfootnote{%
\textit{(1)}~densitates %
\textit{(2)}~densitas aeris post alterationem %
\textit{(3)}~densitates aeris %
\textit{(4)}~Jam vero%
~\textit{L}}}%
\pend
\normalsize
\vspace{1em}
\pstart
Jam %
vero %
\edlabel{LH_37_04_169-170_26b}%
calores\indexS{calor} aeris\indexS{aer} sunt in ratione composita %
ponderum\indexS{ratio ponderum} coercentium et voluminum\indexS{ratio voluminum} ejusdem%
\edtext{aeris\indexS{aer} %
\edlabel{LH_37_04_169-170_22a}%
(per prop. ... cap. de causa et Effectu).}%
{\lemma{prop. ... Effectu}
\Cfootnote{Auslassungspunkte von Leibniz. \textsc{G.~W. Leibniz}, \textit{Dynamica de potentia et legibus}, pars\,II, sectio\,I, prop.\,25; (%
\textit{LMG}\cite{01043}\,VI, S.\,450\cite{01354}; %
\textit{CFP}~I, S.\,331\:\lbrack A\rbrack. S.\,572\:\lbrack B\rbrack, II.2, S.\,399\textendash401\:\lbrack C\rbrack\cite{01482}).}%
{\lemma{aeris}%
\Bfootnote{%
\textit{(1)}~(per praecedentem). %
\textit{(2)}~(per \lbrack...\rbrack\ Effectu).%
~\textit{L}}}}
\edlabel{LH_37_04_169-170_22b}%
Ergo calor aeris\indexS{calor aeris} post alterationem %
est ad calorem\indexS{calor} aeris\indexS{calor aeris} %
\edtext{priorem in ratione }%
{\lemma{priorem}%
\Bfootnote{%
\textit{(1)}~ut %
\textit{(a)}~in ra %
\textit{(b)}~r %
\textit{(c)}~ut \textit{A} %
\textit{(2)}~in ratione%
~\textit{L}}}%
\edtext{composita \textit{B} ad \textit{A} et \textit{M} ad \textit{L}}%
{\lemma{composita}%
\Bfootnote{%
\textit{(1)}~\textit{A} ad \textit{B} et \textit{L} ad \textit{M} %
\textit{(2)}~\textit{B} \lbrack...\rbrack\ \textit{L}%
~\textit{L}}}
seu ut factum ex %
\edtext{numero \textit{B} in \textit{M}}%
{\lemma{numero}%
\Bfootnote{%
\textit{(1)}~\textit{A} in \textit{L} %
\textit{(2)}~\textit{B} in \textit{M}%
~\textit{L}}}
ad %
factum ex %
\edtext{numero \textit{A} in \textit{L}.}%
{\lemma{numero}%
\Bfootnote{%
\textit{(1)}~\textit{B} in \textit{M} \bars{multiplicatio~\textit{erg. u. gestr.}}, %
\textit{(2)}~\textit{A} in \textit{L},%
~\textit{L}}}
\edlabel{LH_37_04_169-170_11a}%
\pend%
\vspace{1em}
\pstart
\noindent
\lbrack\textit{Nachfolgend klein gedruckter Text in L gestrichen:}\rbrack%
\pend
\vspace{0.25em}
\pstart
\footnotesize%
\edtext{%
Factus autem ex \textit{B} in \textit{M} est factus ex \textit{A} plus (vel minus) bis \textit{N} multiplicato per \textit{L} plus (vel minus) \textit{N}.%
}{%
\xxref{LH_37_04_169-170_11a}{LH_37_04_169-170_11b}%
\lemma{\textit{L}}%
\Bfootnote{%
\textit{(1)}~,~seu %
\textit{(a)}~ut
\textit{(aa)} \textit{A} %
\textit{(bb)}~\textit{L} est ad %
\textit{(b)}~ut \textit{A} in \textit{L} est ad \bars{\textit{A}~\textit{streicht Hrsg.}} plus (vel minus) \textit{N} %
\textit{(2)}~. Factus \lbrack...\rbrack\ \textit{N}.
~\textit{L}}}
\edlabel{LH_37_04_169-170_11b}%
Qui est \textit{A} in %
\edtext{\textit{L}, plus (vel minus) \textit{A} in \textit{N}, plus (vel minus) bis \textit{L} in \textit{N}}%
{\lemma{\textit{L},}%
\Bfootnote{%
\textit{(1)}~plus \textit{A} in \textit{N}, plus bis \textit{L} in \textit{N}, plus %
\textit{(2)}~plus \lbrack...\rbrack\ \textit{N},%
~\textit{L}}}
plus bis \textit{N} %
\edlabel{LH_37_04_169-170_12a}%
quadrat. %
\edtext{Itaque ut gradum caloris\indexS{gradus caloris}%
}{%
\xxref{LH_37_04_169-170_12a}{LH_37_04_169-170_12b}%
\lemma{quadrat.}%
\Bfootnote{%
\textit{(1)}~Et %
\textit{(2)}~Itaque \bars{per calculum facilem~\textit{erg. u. gestr.}} %
\textit{(3)}~Si calor ordinarius ponatur esse ut unitas, fiet calor novus %
\textit{(3)}~Itaque %
\textit{(4)}~dato \textit{N} numero partium %
\textit{(5)}~Itaque ut gradum caloris%
~\textit{L}}}
\edlabel{LH_37_04_169-170_12b}%
ascribamus ad quodlibet punctum thermometri\indexS{thermometrum} %
\textit{P} (vel \textit{\greek{f}}) dato numero partium aequalium %
\edtext{inter \textit{Q} et \textit{P}}%
{\lemma{inter}%
\Bfootnote{%
\textit{(1)}~\textit{P\greek{f}} %
\textit{(2)}~\textit{Q} et \textit{P}%
~\textit{L}}}
interceptarum nempe \textit{N}, et calore\indexS{calor} %
\edtext{ordinario Hydrargyrum\indexS{hydrargyrum}}{\lemma{}%
\Bfootnote{%
ordinario \bars{enim~\textit{gestr.}} Hydrargyrum\indexS{hydrargyrum}%
~\textit{ L}}}
sustinente in \textit{Q} posito graduum\indexS{gradus} 100 %
et \textit{A} numero partium Hydrargyri\indexS{hydrargyrum} ordinarii posito 1000, et %
fiet numerus gradus caloris\indexS{gradus caloris} novi.%
\edlabel{LH_37_04_169-170_14a}%
\pend
\pstart
\footnotesize
\edtext{Unde si ponamus calorem\indexS{calor} ordinarium, (nempe %
\edlabel{LH_37_04_169-170_14b}
}{%
\xxref{LH_37_04_169-170_14a}{LH_37_04_169-170_14b}%
\lemma{novi.}%
\Bfootnote{%
\textit{(1)}~Datur autem %
\textit{(2)}~Posito %
\textit{(3)}~Datur %
\textit{(a)}~autem %
\textit{(b)}~vero et \textit{B}, seu numerus partium quas continet \textit{FG}, qui est aequalis ut \textit{A} plus (vel minus) bis \textit{N}, %
\textit{(4)}~\bars{Unde si~\textit{erg.~L}} ponamus %
\textit{(a)}~autem %
\textit{(b)}~jam %
\textit{(c)}~calorem ordinarium %
\textit{(aa)}~ut %
\textit{(bb)}~(nempe%
~\textit{L}}}%
Hydrargyro\indexS{hydrargyrum} existente %
ad altitudinem communem, seu %
\edtext{in \textit{P}),}%
{\lemma{in}%
\Bfootnote{%
\textit{(1)}~\textit{Q} %
\textit{(2)}~\textit{P})%
~\textit{L}}}
esse %
graduum\indexS{gradus} %
\edtext{100. Calor\indexS{calor} novus erit}{\lemma{}%
\Bfootnote{%
100. \bars{ergo~\textit{gestr.}} Calor novus %
\textit{(1)}~est %
\textit{(2)}~erit%
~\textit{ L}}}
graduum\indexS{gradus} quorum numerus fit dividendo numerum factum %
ex multiplicatione 
\edtext{\textit{B} et \textit{M} invicem;}%
{\lemma{\textit{B}}%
\Bfootnote{%
\textit{(1)}~in \textit{M} %
\textit{(2)}~per \textit{M}, %
\textit{(3)}~in %
\textit{(4)}~et \textit{M} %
\textit{(a)}~in se %
\textit{(b)}~invicem;%
~\textit{L}}}
per decuplum ipsius \textit{L}. Is ergo numerus %
ascribetur puncto cui libuerit%
\edtext{\textit{F} supra}{\lemma{\textit{}} %
\Bfootnote{%
\textit{F} \bars{vel \textit{G}~\textit{gestr.}} supra%
~\textit{ L}}}
\textit{P}, %
\edtext{vel puncto \textit{\greek{f}}}%
{\lemma{vel}%
\Bfootnote{%
\textit{(1)}~\textit{F} %
\textit{(2)}~puncto \textit{\greek{f}}%
~\textit{L}}}
infra \textit{P}; hoc tantum observato, ut quando %
diximus: plus vel minus serviat plus pro \textit{F}, et minus pro %
\edtext{\textit{\greek{f}}. Ex his vicissim facili calculo habetur praxis talis, si quis non dato spatio \textit{N} quaerat gradum caloris\indexS{gradus caloris}, sed contra dato gradu caloris\indexS{gradus caloris}, quaerat \textit{N} magnitudinem ipsius \textit{PF} (vel \textit{P\greek{f}}). %
\edlabel{LH_37_04_169-170_16a}%
Nempe assumat %
\edtext{numerum \textit{Y}}{%
\lemma{numerum \textit{Y}}%
\Cfootnote{Der Buchstabe \textit{Y} bezeichnet hier wie an weiteren gestrichenen Stellen dieselbe Größe, für die im gültigne Text des Stückes \textit{V} steht. Die Bezeichnung \textit{Y} geht auf N.~\ref{60326} zurück, wo Leibniz auf \refpassage{LH_37_05_108_01a}{LH_37_05_108_01b} anmerkt, diese Bezeichnung durch den Buchstaben \textit{V} in der fertigen Fassung (N.~\ref{60270}) ersetzt zu haben.}}
qui aequivaleat quartae parti ipsius \textit{A} et dimidiae ipsius \textit{L} simul sumtis \edlabel{LH_37_04_169-170_16b} deinde et numerum $\varepsilon$ qui sit differentia gr. inter numeros graduum caloris\indexS{gradus caloris} novi, et caloris\indexS{calor} ordinarii, inde ex numero qui sit quadratum numeri \textit{Y}, plus vel minus facto ex numeris, \textit{A}, \textit{L}, $\varepsilon$, in se invicem multiplicatis extracta radix quadratica, plus vel minus \textit{Y}, erit \textit{N}. Unde et Tabulam\indexS{tabula} pro 100 gradibus\indexS{gradus} condere licebit. Sed si quis malit statui ordinario caloris\indexS{calor} seu medio quem sensus noster nec calidum nec frigidum iudicat ascribi non 100, sed 0 
tunc pro gradibus\indexS{gradus} 110, 120, etc. scribentur gradus caloris\indexS{gradus caloris} 10, 20, etc. At loco calorum\indexS{calor} imminutorum 90, 80 etc. poterunt scribi gradus frigoris\indexS{gradus frigoris} 10, 20 etc., et alterutro modo habebimus thermometrum\indexS{thermometrum} in gradus\indexS{gradus}}%
{\lemma{\textit{\greek{f}}}%
\Bfootnote{%
\textit{(1)}~et condi poterit tabula pro gradibus 100 et habebimus thermometrum\indexS{thermometrum} in gradus accurate divisum, ubi revera numeri minorum calorum significabunt frigus\indexS{frigus}: %
\textit{(a)}~quos e %
\textit{(b)}~et pro gradibus caloris infra \textit{P} in \textit{\greek{f}} poterunt ascribi gradus frigoris. Sumendo pro frigore numerus reciproce proportionales caloris, ut si calor \bars{ordinarius (seu potius status medius per sensu nostro inter frigus et calorem \textit{Y}) \textit{(1)}~esset inv \textit{(2)}~quem \textit{(a)}~concipimus \textit{(b)}~ponimus graduum 100 factus esset~\textit{erg.}} esset graduum %
\textit{(aa)}~80 %
\textit{(bb)}~90 %
\textit{(2)}~Ex \lbrack...\rbrack\ talis, \bars{et generaliter~\textit{gestr.}} si quis non dato spatio \bars{\textit{N}~\textit{erg.}} \lbrack...\rbrack\ contra %
\textit{(a)}~quod %
\textit{(aa)} si quis %
\textit{(bb)} autem %
\textit{(b)}~dato gradu \lbrack...\rbrack\ ipsius %
\textit{(aa)}~\textit{LP} %
\textit{(bb)}~\textit{Q} %
\textit{(cc)}~\textit{PF} (vel \textit{P\greek{f}}). \bars{Nempe~\textit{erg.}} assumat \lbrack...\rbrack\ ipsius \textit{L} \bars{simul sumtis~\textit{erg.}} deinde et \lbrack...\rbrack\ medio %
\textit{(aaa)}~pro sensu nostro %
\textit{(bbb)}~quem \lbrack...\rbrack\ iudicat \lbrack...\rbrack\ etc. %
\textit{(aaaa)}~ascribe %
\textit{(bbbb)}~scribetur 10, 20 etc. %
\textit{(cccc)}~scribentur \lbrack...\rbrack\ etc. %
\textit{(aaaaa)}~et habebimus thermometrum in gradus %
\textit{(aaaaa-a)}~accurate divisum %
\textit{(bbbbb-b)}~aequales %
\textit{(bbbbb)}~et \lbrack...\rbrack\ gradus%
~\textit{L}}}
aequales per spatia
\edtext{licet}{\lemma{}%
\Bfootnote{licet~\textit{erg.~L}}}
inaequalia accurate divisuri %
parato autem semel hoc Thermometro\indexS{thermometrum}, aliae %
\edtext{alterius constructionis,}%
{\lemma{alterius}%
\Bfootnote{%
\textit{(1)}~molis %
\textit{(2)}~magnitudini %
\textit{(3)}~constructionis%
~\textit{L}}}
et minoris molis, ope ejus %
accurate dividi possunt quae facienda nobis proposueramus.
\pend
\vspace{0.5em}
\normalsize
\pstart
Itaque posito calorem medium aeris\indexS{calor aeris} seu ordinarium, (nempe %
cum Hydrargyri\indexS{hydrargyrum} altitudo est \textit{PQ}) 
\edtext{esse ut unitatem.}%
{\lemma{esse}%
\Bfootnote{%
\textit{(1)}~graduum 100 %
\textit{(2)}~ut unitatem. %
\textit{(a)}~Itaque %
\textit{(b)}~Hinc%
~\textit{L}}}
Hinc si \textit{B} multiplices per %
\edtext{\textit{M} et productum denique dividas}%
{\lemma{\textit{M}}%
\Bfootnote{%
\textit{(1)}~, et quod provenit per 100, %
\textit{(2)}~et productum denique dividas%
~\textit{L}}}
per factum ex \textit{A} multiplicato %
per \textit{L}, quotiens proveniens erit gradus caloris\indexS{gradus caloris} %
novi quaesitus. Itaque sic data alteratione Hydrargyri\indexS{hydrargyrum} %
nempe \textit{N}, id est \textit{PF} (vel \textit{P\greek{f}}) elevatione eius (vel depressione) %
datur gradus caloris\indexS{gradus caloris}, sed %
\edtext{nobis jam potius opus est ad gradus caloris\indexS{gradus caloris} et frigoris in thermometro\indexS{thermometrum} designandos dato}%
{\lemma{nobis}%
\Bfootnote{%
\textit{(1)}~\bars{jam~\textit{streicht Hrsg.}} vicissim opus est ad divisionem Ther %
\textit{(2)}~jam potius opus est %
\textit{(a)}~dato \textit{N} %
\textit{(b)}~ad \lbrack...\rbrack\ dato%
~\textit{L}}}
caloris\indexS{calor} gradu\indexS{gradus caloris} invenire \textit{N}, seu \textit{PF} vel \textit{P\greek{f}}, id est %
punctum \textit{F}, vel \textit{\greek{f}}, cui %
\edtext{ille gradus\indexS{gradus} debet ascribi.}%
{\lemma{ille}%
\Bfootnote{%
\textit{(1)}~calor debet ascribi %
\textit{(2)}~gradus debet ascribi%
~\textit{L}}}
Quod facili %
calculo praestatur hoc %
\edtext{modo $\pm B \mp A=N$ ex hypothesi, ergo $\displaystyle B=\pm N+A$. Jam $\pm M \mp L=2N$}%
{\lemma{modo}%
\Bfootnote{%
\textit{(1)}~ %
\textit{(a)}~$+\,A\,-$ %
\textit{(b)}~$\pm A \mp B=N$, ex hypothesi prima %
\textit{(2)}~, \bars{primo \textit{streicht Hrsg.}} %
\textit{(a)}~$\pm A$ %
\textit{(b)}~\bars{(1) \textit{streicht Hrsg.}} $\pm B$ \bracks{...}, ergo \bars{(secunda)~\textit{gestr.}} $B=\pm N+A$. Jam \bars{(tertio)~\textit{gestr.}} $\pm M \mp L=2N$%
~\textit{L}}}
\edtext{(ex supra ostensis). }{%
\lemma{ex supra ostensis}%
\Cfootnote{Siehe in vorliegendem Stück \refpassage{LH_37_04_169-170_18a}{LH_37_04_169-170_18b}.}}%
\rule[0mm]{0pt}{5.5mm}%
Ergo $\displaystyle M=\pm2N+L$. %
\edlabel{LH_37_04_169-170_19a}%
Ac proinde %
$\displaystyle BM=2NN\pm\overline{2A+L}\;N+AL$ jam ex %
\edtext{supra inventis }{%
\lemma{supra inventis}%
\Cfootnote{Siehe in N.~\ref{60326} \refpassage{LH_37_05_108_02a}{LH_37_05_108_02b}}}
erat %
\edtext{$BM=cAL$.}{\lemma{%
}\Afootnote{%
\textit{Am Rand}:
$c:100\squaredots BM:AL$\hspace{3mm} $c\cdot AL=100\cdot BM$}}
Fit ergo $\displaystyle NN\pm\frac{2A+L}{2}\ N=\frac{1}{2}\ cAL-\frac{1}{2}\ AL$. %
\edlabel{LH_37_04_169-170_19b}%
\edlabel{LH_37_04_169-170_17a}%
Sit \rule[0mm]{0pt}{5.5mm}%
 $\displaystyle \frac{2A+L}{4}=V$, \edlabel{LH_37_04_169-170_17b}%
et dimidia differentia inter \textit{c} et 1 seu %
inter gradum\indexS{gradus} ordinarium et novum sit\rule[0mm]{0pt}{3mm} %
$\varepsilon$, erit $c-1\ \text{aequ.}+2\:\!\varepsilon$ %
in casu caloris\indexS{calor} aucti, et $c-1\ \text{aequ.}\,-2\:\!\varepsilon$ in casu caloris\indexS{calor} %
diminuti. Ergo alterutro modo (ex natura signorum %
ambiguorum hic assumtorum\lbrack)\rbrack, fit $c-1\ \text{aequ.}\,\pm2\:\!\varepsilon$, ac %
proinde fit $NN\pm2VN=\,\pm\,\varepsilon AL$. Unde haec nascitur operatio %
posito gradum caloris\indexS{gradus caloris} ordinarium seu medium repraesentari %
per unitatem, quem dato gradu aucti vel diminuti caloris\indexS{gradus caloris} quaeritur %
longitudo rectae \textit{PF} vel \textit{P\greek{f}}, seu spatium quo hydrargyrum\indexS{hydrargyrum} calore\indexS{calor} %
aucto assurgere supra \textit{P}, vel diminuto %
descendere infra \textit{P} oportet, ut hic propositus calor\indexS{calor} a thermometro\indexS{thermometrum} %
indicetur. %
Volumen aeris\indexS{volumen aeris} ordinarii, seu spatium %
\edtext{\textit{ZXVT}\ 
\href{\Scan{LH_37_04_169r}}%
{\lbrack169\:r\textsuperscript{o}\rbrack}\
expressimus numero \textit{A} (notante quot verbi gratia pollices}%
{\lemma{\textit{ZXVT}}%
\Bfootnote{%
\textit{(1)}~vocavimus \bars{\textit{A}~\textit{streicht Hrsg.}}, et volumen ordinarium Hydrargyri s %
\textit{(2)}~\bars{expressimus~\textit{streicht Hrsg.}} %
\textit{(a)}~numero partem \ 
\href{\Scan{LH_37_04_169r}}%
{\lbrack169~r\textsuperscript{o}\rbrack}\
suspensi, seu spatium \textit{PQ} vocavimus \textit{L} %
\textit{(b)}~sus %
\textit{(c)}~numerum partium (verbi gratia pollicum %
\textit{(3)}~expressimus numero \textit{A} %
\textit{(a)}~(exprimente %
\textit{(b)}~(notante quot verbi gratia pollices%
~\textit{L}}}
cubicos hoc spatium contineat, quemadmodum prius aquam %
experiri licuit) et volumen ordinarium Hydrargyri\indexS{hydrargyrum} %
suspensi, seu spatium \textit{PQ}, expressimus numero \textit{L}.  %
\edtext{Jam dimidiam}%
{\lemma{Jam}%
\Bfootnote{%
\textit{(1)}~quartam %
\textit{(2)}~dimidiam%
~\textit{L}}}
partem ipsius \textit{A}, auctam quarta %
parte ipsius %
\edtext{\textit{L} designemus}%
{\lemma{\textit{L}}%
\Bfootnote{%
\textit{(1)}~vocemus %
\textit{(2)}~designemus%
~\textit{L}}}
numero \textit{V}, et %
denique differentiam dimidiam inter calorem\indexS{calor} medium %
per unitatem repraesentatum et gradum\indexS{gradus} novum %
caloris\indexS{calor} aucti vel minuti vocemus $\varepsilon$. %
\edtext{Jam quadrato}{\lemma{Jam}%
\Bfootnote{%
\textit{(1)}~ad quadrata %
\textit{(2)}~quadrato%
~\textit{L}}}
ipsius \textit{V} addatur in casu caloris\indexS{calor} aucti %
detrahatur in casu caloris\indexS{calor} imminuti factum ex %
tribus numeris \textit{A}, \textit{L}, $\varepsilon$; et ex numero hac additione %
vel subtractione proveniente extrahatur radix %
\edtext{quadrata, et differentia inter}{\lemma{quadrata,}%
\Bfootnote{%
\textit{(1)}~cui in casu aucti %
\textit{(2)}~et differentia inter %
~\textit{L}}} %
hanc radicem et numerum \textit{V}, erit numerus \textit{N} %
quaesitus, indicans quantitatem spatii \textit{PF} vel \textit{P\greek{f}}, %
adeoque punctum \textit{F} vel %
\edlabel{LH_37_04_169-170_20a}%
\edtext{%
\edtext{\textit{\greek{f}}. Et exempli causa}{\lemma{\textit{\greek{f}}}%
\Bfootnote{%
\textit{(1)}~Quod si ponimus inter gradus %
\textit{(2)}~Et exempli causa%
~\textit{L}}}
sit \textit{A} 100, \textit{L} 60, \textit{c} 2 (seu calor\indexS{calor} ordinarius %
\edtext{sit}{\lemma{}\Bfootnote{sit \textit{ erg.} \textit{ L}}} %
duplicatus) erit \textit{V}\:65, %
$\rule[0mm]{0pt}{6mm}\varepsilon\:\displaystyle\frac{1}{2}$. Quadratum a 65 est 4225 cui addatur factum ex 100, 60, 
$\rule[0mm]{0pt}{2mm}\displaystyle\frac{1}{2}$, id est 3000, %
fit 7225, cujus radix quadratica est %
\edtext{85, unde}%
{\lemma{85,}%
\Bfootnote{%
\textit{(1)}~unde (\textit{nicht gestr.}) %
\textit{(2)}~seu %
\textit{(3)}~unde%
~\textit{L}}}
si auferas 65, restat 20. %
}{%
\lemma{Et exempli causa}%
\Cfootnote{%
Die Rechungen hierzu macht Leibniz in N.\:\ref{60326}, \refpassage{LH_37_05_108_05a}{LH_37_05_108_05b} und \refpassage{LH_37_05_108_04a}{LH_37_05_108_04b}.
}}%
\edlabel{LH_37_04_169-170_20b}%
\rule[0mm]{0pt}{4mm}%
Ergo in casu propositi thermometri\indexS{thermometrum} si calor\indexS{calor} sit 2 seu duplicatus %
\edtext{tunc 20,}{\lemma{}%
\Bfootnote{tunc 20~\textit{erg.~L}}}
erit %
\textit{PF} ascensus Mercurii\indexS{mercurius} ultra \textit{QP} altitudinem ordinariam, %
quod si velimus gradus frigoris\indexS{gradus frigoris} non considerare ut gradus %
\edtext{60}{\lemma{}%
\Bfootnote{60~\textit{erg.~L}}}
medii caloris\indexS{gradus caloris} diminuti, sed ut gradus\indexS{gradus} propriae %
qualitatis, 
\edtext{tunc}{\lemma{}%
\Bfootnote{tunc~\textit{erg.~L}}}
gradus medii status %
exprimetur per 0, et gradus caloris\indexS{gradus caloris} %
et frigoris sic %
\edtext{designandi, erunt differentiae inter paulo ante inventos et unitatem. }%
{\lemma{designandi,}%
\Bfootnote{%
\textit{(1)}~ex %
\textit{(2)}~habebantur, %
\textit{(3)}~si a paulo ante inventis detrahatur unitas, ubi residuum majus nihilo dabit  g %
\textit{(4)}~si eorum qui paulo ante inventi sunt differentia ab unitate quaeratur, quae erit gradus caloris si excessus differentiam pr\textlangle odit\textrangle\ gradum antea inventum %
\textit{(5)}~erunt differentiae %
\textit{(a)}~paulo ante %
\textit{(b)}~inter \lbrack...\rbrack\ unitatem.%
~\textit{L}}}%
Et quidem si gradus\indexS{gradus} paulo ante inventi majores unitate fuerunt, prodeunt gradus caloris\indexS{gradus caloris} sin minores habemus gradus frigoris\indexS{gradus frigoris}. Quorum illi supra \textit{P} ut in \textit{F}, hi infra ut in \textit{\greek{f}}, designantur, potest etiam hac methodo tabula\indexS{tabula} condi ex qua numeri in %
\edtext{instrumentum\indexS{instrumentum} transferentur.}%
{\lemma{instrumentum}%
\Bfootnote{%
\textit{(1)}~transferri possunt. %
\textit{(2)}~transferentur.%
~\textit{L}}}
Parato jam semel hoc Thermometro\indexS{thermometrum} %
praescripto, tanquam mensura caeterorum %
\edtext{et sui ipsius regula,}{\lemma{}%
\Bfootnote{et sui ipsius regula,~\textit{erg.~L}}}
alia alterius constructionis cujuscunque et minoris operae, cum hoc collata %
similiter in gradus caloris\indexS{gradus caloris} (aequales) %
accurate dividi poterunt. Quod nobis %
faciendum proposueramus.  %
\pend%
\vspace{1em}
\pstart
\noindent
\lbrack\textit{Am Rand von Bl.~\href{\Scan{LH_37_04_170v}}{170\:v\textsuperscript{o}}:}\rbrack%
\pend
\vspace{0.5em}
\pstart\noindent
Sit %
$\displaystyle L+\frac{1}{2}A=R$ %
fiet%
\edlabel{LH_37_04_169-170_06a}%
\edtext{}{%
\xxref{LH_37_04_169-170_06a}{LH_37_04_169-170_06b}%
\lemma{fiet}%
\Bfootnote{%
\textit{(1)}~$\frac{1}{2}A^2+RN+NN$ %
\textit{(2)}~$LR$ %
\textit{(3)}~$\MP RN+NN$ 
~\textit{L}}}
$\displaystyle \MP RN+NN\edlabel{LH_37_04_169-170_06b}\edlabel{LH_37_04_169-170_23a}-\frac{1}{2}AL+\frac{1}{2}AL\cdot c\edlabel{LH_37_04_169-170_23b}$ %
seu%
\edtext{}{%
\xxref{LH_37_04_169-170_23a}{LH_37_04_169-170_23b}%
\lemma{}%
\Bfootnote{%
$-\frac{1}{2}AL$ \bars{$+\frac{1}{4}RR$ \textit{streicht Hrsg.}}$+\frac{1}{2}AL\cdot c$~\textit{L}}}
$\displaystyle \edlabel{LH_37_04_169-170_25a}N=\frac{1}{2}\sqrt{R^2+2AL\edlabel{LH_37_04_169-170_24a}\overline{[1-c]}\edlabel{LH_37_04_169-170_24b}}\;\PM\frac{1}{2}R\edlabel{LH_37_04_169-170_25b}$.%
\edtext{}{%
\xxref{LH_37_04_169-170_25a}{LH_37_04_169-170_25b}%
\lemma{$N=$}%
\Bfootnote{%
\textit{(1)}~$\MP R$
\textit{(2)}~$\PM R$
\textit{(3)}~$\frac{1}{2}\sqrt{R^2+2AL\overline{[1-c]}}\;\PM\frac{1}{2}R$~\textit{L}}}
\edtext{}{%
\xxref{LH_37_04_169-170_24a}{LH_37_04_169-170_24b}%
\lemma{}%
\Bfootnote{%
$-1+c$ 
\textit{L ändert Hrsg.}}}
Sit\,$\displaystyle\frac{1}{2}\;\!R=Y$ et differentia %
\edtext{dimidia}{\lemma{}\Bfootnote{dimidia \textit{erg.~L}}}
\edtext{inter 1 et \textit{c}}{%
\lemma{inter}\Bfootnote{%
\textit{(1)}~100 et \textit{c} %
\textit{(2)}~1 et \textit{c} %
~\textit{L}}}
seu inter gradum caloris\indexS{gradus caloris} novum et gradum caloris\indexS{gradus caloris} ordinarium %
repraesentatum per unitatem vocetur %
\edtext{$\varepsilon$, sin sit $\varepsilon=+1-c$ in casu quo $\PM$ significat $+$ et $-1+c$ in casu quo $\MP$ significat $+$.}{%
\lemma{$\varepsilon$,}%
\Bfootnote{%
\textit{(1)}~fit $N=\surd{YY\PM}\; \varepsilon\;\PM$ %
\textit{(a)}~$\varepsilon\varepsilon$ %
\textit{(b)}~$\varepsilon$ %
\textit{(2)}~Ergo $c\PM \varepsilon$ %
\textit{(3)}~Ergo $-1+c\;\PM$ %
\textit{(4)}~Ergo %
\textit{(a)}~$-c$
\textit{(b)}~$+c-1=+\:\!\varepsilon$ in casu quo $\PM=+$. Ergo $c-1=\MP\:\!\varepsilon$ et fit $\displaystyle N=\sqrt{YY\MP \varepsilon\cdot AL}\PM Y$. %
\textit{(5)}~sin sit \lbrack...\rbrack\ significat $+$.%
~\textit{L}}}
\pend
\vspace{2em}
\pstart
\noindent
\lbrack\textit{Gestrichene Rechnungen am Rand von Bl.~\href{\Scan{LH_37_04_170v}}{170\:v\textsuperscript{o}}:}\rbrack%
\pend
\vspace{0.5em}%
\footnotesize
\pstart\noindent
$\displaystyle \langle c\rangle: 1\squaredots BM:AL$. %
Ergo %
$\displaystyle c\cdot AL=BM$, jam $A-B\PM A\MP B=N=\PM L\MP M$. %
Ergo %
$\displaystyle B=A\MP N$. %
Et %
$\displaystyle M=L\MP2N$. %
Ergo %
$\displaystyle BM=AL\MP AN+2NN=AL\,\cdot\,c$. %
Seu %
$\displaystyle NN\MP \overline{L+\displaystyle\frac{1}{2}A}\ N+\displaystyle\frac{1}{4}LL+\displaystyle\frac{1}{4}AL+\displaystyle\frac{1}{4}AA=\displaystyle\frac{1}{4}LL+\displaystyle\frac{1}{4}AL+\displaystyle\frac{1}{4}AA+\displaystyle\frac{1}{2}ALc:100$ %
et %
$\displaystyle N=\frac{1}{2}\protect\surd{\llcorner LL+AA+\overline{L+2c:100}\ NL\lrcorner}\PM L+\frac{1}{2}A$ %
\pend
\normalsize
